\section{The experimental problem}\label{sec:experiment}
The mother problem is to pass from a physically executable experiment to a quantitative law for the predictions that can actually be retained. There are two distinct obstructions. A set of formally reachable predictions need not receive appreciable mass under the preparation and exploration being used. Moreover, a good codebook at a checkpoint need not admit a good recursive implementation. The first obstruction is probabilistic; the second is causal. Treating either as a dimension count loses the mechanism of the problem.

Our main theorem connects these obstructions by weighting the propagation of quantization errors with the same acquisition probabilities that enter the lower bound. Inactive steps are especially important. When an update preserves a retained representative exactly, rounding need not be repeated. Charging an error at every physical tick would produce a spurious inverse acquisition probability in a system that simply waits. The theorem retains the actual clock while localizing error injection to the updates that require it.

The principal mechanism has two parts. A positive-kernel resolvent computes the energy of past rounding innovations from one-step experimental data. Its Lyapunov hypothesis permits expansion, rather than assuming that every informative update contracts. A separate continuation-separation argument decides when additional interface states are genuinely mandatory. This prevents a Cartesian-product implementation bound from masquerading as a memory lower bound. These mechanisms yield, respectively, a general causal transfer theorem and a common-task robust minimax theorem with growing total-state overhead.

Predictive states and their minimality have classical precedents in causal-state theory and controlled predictive representations \cite{shalizi,psr}. Conditional projection and probability quantization are also classical \cite{graf}, as are finite filter approximation \cite{pages,sellami,feuer} and comparison of experiments \cite{blackwell,lecam,torgersen}. Positive-operator moment methods are related to mean-square stability of Markov jump systems \cite{costa}. The contribution asserted here is the connection of these mechanisms to actual acquisition mass, forced causal approximation, and a precisely specified resource cut. The theorem-level distinctions are given in Section~\ref{sec:scope}; no priority is claimed for the component theories.

\subsection{Raw data, strategies, and preparation}
A discrete-time experiment consists of standard Borel spaces $X_t,\mathcal A_t,Y_{t+1}$, a probability preparation $\mu$ on $X_0$, and nonnegative normalized kernels
\begin{equation}\label{eq:raw}
 K_t(dx',dy,dc\mid x,a).
\end{equation}
The visible history $h_t$ records actions, reports, and the published nonnegative resource increments $c$. These increments include failed preparations and physical elapsed time. Positivity means nonnegativity and normalization, not a positive lower bound for every report. Deterministic transitions are permitted. An unknown physical parameter, when present, is included in the hidden state with a specified prior for the Bayesian assertions below. Parameter-uniform experiment comparison will be stated separately.

A legal strategy is a stochastic kernel $\alpha_t(da\mid h_t)$ supported on the declared available actions. All randomness used by the apparatus or algorithm is independent of the hidden preparation unless generated by \eqref{eq:raw}. A stopping rule is a stopping time for the visible filtration. The main risk theorem fixes a deterministic raw-acquisition horizon $N$ and an exploration strategy $\beta$ independently of the encoder. It does not optimize $\beta$. At most $N$ actual calls, including unsuccessful calls, are made. A stopped comparison is always explicitly bounded by $N$.

Successive kernel integration gives the path law $P^\alpha$. The complete conditional belief $\pi_t(h_t)$ on the hidden state has a measurable version. We specify versions on a declared admissible history domain; identities on impossible reports are not inferred by division by zero.

\subsection{Executable future tests and the quotient}
A future test is a legal finite continuation followed by a bounded measurable function of its actual reports. Hidden-state functions qualify only when a terminal measurement producing them is part of the experiment. Tests include determining report events and the continuation interface: legal actions, remaining budget, and every clock value affecting future experiments. They are closed under legal one-step continuation, finite randomization, and the bounded monotone limits needed to determine the designated laws.

For an admissible prepared belief $b$, let $\mathsf e_t(b)$ be the evaluations of a countable determining family of these tests, together with the continuation interface. Equality of $\mathsf e_t$ is \emph{predictive equivalence}. The family determines the entire designated continuation law, not merely the finite scoring menu introduced below. Here and throughout, a statement about an unavailable action is not a continuation assertion.

We impose the following realization certificate, verified from the instruments. The admissible belief domain $\mathcal B_t$ and the action domain are compact metric. The determining evaluations are continuous on $\mathcal B_t$. For a continuous report component there is a compact metric report domain $Y$, a finite dominating measure $\lambda$ with full support on $Y$, and an instrument density $T_{t,a,y}$. For every determining next-time test $f$, the functions
\begin{equation}\label{eq:density-pullback}
 n_f(b,a,y)=bT_{t,a,y}f,\qquad g(b,a,y)=bT_{t,a,y}1
\end{equation}
are jointly continuous. The functions $b\mapsto\int_E n_f(b,a,y)\lambda(dy)$, for report events $E$ in a determining algebra, are evaluations of executable one-step continuations, hence are constant on predictive fibres. The same is required for $f=1$. This is an \emph{event-instrument closure}, not the assertion that a report singleton has positive probability. The normalized belief update $U_{t,a,y}$ preserves the declared domain and is jointly continuous where $g>0$. If it is used at $g=0$, a continuous, predictive-fibre-compatible extension is supplied separately. Finite report components use the corresponding atomic instruments. Deterministic reveal components may instead be handled by an explicitly verified continuous quotient update, as in Section~\ref{sec:sensor}.

\begin{lemma}[Measurable factorization with completed histories]\label{lem:factor}
Let $R:\Omega\to E$ and $S:\Omega\to S_0$ be random variables, where $E$ is standard Borel and $S_0$ is a nonempty Borel subset of a countable product of compact intervals. If every coordinate of $S$ is measurable with respect to the completion of $\sigma(R)$ under $P$, there is a Borel map $f:E\to S_0$ such that $S=f(R)$ almost surely. It is unique $P_R$-almost everywhere among factors of $S$.
\end{lemma}
\begin{proof}
A completed-measurable real random variable has an uncompleted-measurable version: replace the finitely many measurable sets in each dyadic simple approximation by $\sigma(R)$-sets differing only by null sets, and take a pointwise limit on the common full-measure set. Every set in $\sigma(R)$ is $R^{-1}(B)$ for some Borel $B\subset E$, since these inverse images already form a sigma algebra. The same simple approximations therefore give Borel coordinate factors on $E$; assign zero wherever their limits fail to exist. Combine the countably many factors into a product-valued Borel map $f_0$. Since $S_0$ is Borel, $\{e:f_0(e)\in S_0\}$ is Borel and has full $P_R$-measure. Replace its complement by a fixed point of $S_0$. This proves existence. Equality of two factors after composition with $R$ is exactly their equality $P_R$-almost everywhere.
\end{proof}

\begin{proposition}[Predictive realization and pointwise update descent]\label{prop:quotient}
Under the preceding certificate, $S_t=\mathsf e_t(\mathcal B_t)$ is compact metrizable and realizes predictive equivalence. Report laws descend continuously. The normalized update descends continuously on the positive-density domain; a declared continuous fibre-compatible extension descends on its declared extended domain. In particular, equivalent beliefs have equivalent updated beliefs for \emph{every} positive-density report, not just almost every report. Any standard-Borel-valued history statistic through which the determining evaluations and continuation interface are completed-measurable admits a measurable predictive factor almost surely. Setwise sufficiency gives a unique setwise factor on the image, but no additional Borel conclusion for an arbitrary measurable quotient.
\end{proposition}
\begin{proof}
Rescale the determining evaluations to $[-1,1]$. Their continuous map into a countable product has compact metrizable image. The determining property identifies its fibres with equality of the designated continuation laws.

Fix equivalent $b,b'$ and an available action $a$. Event-instrument closure gives equal integrals of $n_f(b,a,\cdot)$ and $n_f(b',a,\cdot)$ over a determining algebra. Uniqueness of finite signed measures extends this equality to all Borel events, and equality of densities follows $\lambda$-almost everywhere. A continuous function vanishing almost everywhere under a full-support measure vanishes everywhere: any nonzero value would have a same-sign neighborhood of positive measure. Hence the two numerator functions, and the two denominators, agree for every $y$. Where $g>0$, dividing proves equality of every determining next-time evaluation. The continuation interface is among these evaluations. At a zero denominator only the separately verified extension is used.

A continuous surjection from a compact space onto a Hausdorff space is a quotient map. Its product with a compact action or report domain has the same property. Thus the continuous fibre-constant report evaluations and extended updates descend. Restriction to the saturated open positive-density domain also descends, since that domain is the inverse image of an open set under the descended continuous $g$. Apply Lemma~\ref{lem:factor} to the countable predictive coordinates for measurable minimality. Setwise minimality is the elementary inclusion of sufficient-statistic fibres in predictive fibres.
\end{proof}

This proposition fixes conditional versions rather than inferring them from null events. Finite exact interface components may be included as a disjoint compact union. A noncompact report or belief domain needs a separate Borel realization or a verified localization. No universal smooth realization of an arbitrary predictive equivalence relation is asserted.

\subsection{The common prediction task}
At time $N$, after the label has been retained, an independent index $i\in\{1,\ldots,q\}$ is selected with fixed positive weights $v_i$, $\sum_i v_i=1$. The corresponding legal terminal probe returns $T_i\in[0,1]$. The algorithm reports $a_i\in[0,1]$ and incurs $(T_i-a_i)^2$. Each expectation refers to its own executable probe; a joint counterfactual vector of all probe outcomes is unnecessary. For each possible terminal time $t$, write $T_{t,i}$ for its corresponding executable probe; at the selected horizon abbreviate $T_i=T_{N,i}$. Put
\begin{equation}\label{eq:score}
 p_t(s)=\big(\sqrt{v_i}\,\E[T_{t,i}\mid s]\big)_{i=1}^q,
 \qquad B_N=\sum_i v_i\E\Var(T_i\mid H_N),
 \qquad \nu_N=\law(p_N(s_N)).
\end{equation}
A varying terminal time uses the corresponding declared probe. The finite scoring menu need not determine the entire predictive quotient. Our assumptions below require only its Lipschitz upper bound globally and a separating lower bound on each acquired chart. Unmeasured interface coordinates cannot be silently charged to this task's lower bound.

A checkpoint encoder reads $H_N$ and retains at most $m$ labels. An online encoder has at most $m$ persistent labels at every time, updates from its previous label and the current input/report, and has no readable history tape. A terminal decoder may depend on the public horizon and on the freshly selected probe index, but not on discarded observations. Shared randomness does not carry observation-dependent state. The risks $R_\cp(N,m)$ and $R_\on(N,m)$ are infima of the full expected score over these respective classes, under the fixed exploration $\beta$. Temporary workspace, numerical input access, and program description are recorded separately in Section~\ref{sec:transfer}.

The mathematical label class allows Borel transitions on exact current reports. Its risk is not the risk of a finite-bit processor. For an implementation with finite input precision and arithmetic, we write $R_{\rm alg}$ for its actual risk and retain the errors $u_t$ and $v_N^{\rm num}$ explicitly. The latter is a deterministic upper bound on the terminal score-norm error; a stated $L^2$ bound could be used instead. Appendix~\ref{sec:algorithms} supplies both types of algorithms. A fixed public clock may give the predetermined horizon or protocol phase, but cannot encode observation-dependent past information through adaptive timing. Any such timing channel is retained state and is charged.
